\documentclass[review]{elsarticle}
%\documentclass[twocolumn]{elsarticle}

\usepackage{lineno,hyperref}
\modulolinenumbers[5]



\journal{Icarus}

% 中文支持
\usepackage{xeCJK}
\setCJKmainfont{SimSun}
\setCJKsansfont{Microsoft YaHei}
\setCJKmonofont{SimHei}


% Bibliography and bibfile
\def\aj{AJ}%          % Astronomical Journal
%\def\actaa{Acta Astron.}%          % Acta Astronomica
%\def\araa{ARA\&A}%          % Annual Review of Astron and Astrophys
\def\apj{ApJ}%          % Astrophysical Journal
\def\apjl{ApJ}%          % Astrophysical Journal, Letters
%\def\apjs{ApJS}%          % Astrophysical Journal, Supplement
%\def\ao{Appl.~Opt.}%          % Applied Optics
%\def\apss{Ap\&SS}%          % Astrophysics and Space Science
\def\aap{A\&A}%          % Astronomy and Astrophysics
%\def\aapr{A\&A~Rev.}%          % Astronomy and Astrophysics Reviews
%\def\aaps{A\&AS}%          % Astronomy and Astrophysics, Supplement
%\def\azh{AZh}%          % Astronomicheskii Zhurnal
%\def\baas{BAAS}%          % Bulletin of the AAS
%\def\bac{Bull. astr. Inst. Czechosl.}%          % Bulletin of the Astronomical Institutes of Czechoslovakia 
%\def\caa{Chinese Astron. Astrophys.}%          % Chinese Astronomy and Astrophysics
%\def\cjaa{Chinese J. Astron. Astrophys.}%          % Chinese Journal of Astronomy and Astrophysics
\def\icarus{Icarus}%          % Icarus
%\def\jcap{J. Cosmology Astropart. Phys.}%          % Journal of Cosmology and Astroparticle Physics
%\def\jrasc{JRASC}%          % Journal of the RAS of Canada
%\def\mnras{MNRAS}%          % Monthly Notices of the RAS
%\def\memras{MmRAS}%          % Memoirs of the RAS
%\def\na{New A}%          % New Astronomy
%\def\nar{New A Rev.}%          % New Astronomy Review
%\def\pasa{PASA}%          % Publications of the Astron. Soc. of Australia
%\def\pra{Phys.~Rev.~A}%          % Physical Review A: General Physics
%\def\prb{Phys.~Rev.~B}%          % Physical Review B: Solid State
%\def\prc{Phys.~Rev.~C}%          % Physical Review C
%\def\prd{Phys.~Rev.~D}%          % Physical Review D
%\def\pre{Phys.~Rev.~E}%          % Physical Review E
%\def\prl{Phys.~Rev.~Lett.}%          % Physical Review Letters
%\def\pasp{PASP}%          % Publications of the ASP
%\def\pasj{PASJ}%          % Publications of the ASJ
%\def\qjras{QJRAS}%          % Quarterly Journal of the RAS
%\def\rmxaa{Rev. Mexicana Astron. Astrofis.}%          % Revista Mexicana de Astronomia y Astrofisica
%\def\skytel{S\&T}%          % Sky and Telescope
%\def\solphys{Sol.~Phys.}%          % Solar Physics
%\def\sovast{Soviet~Ast.}%          % Soviet Astronomy
%\def\ssr{Space~Sci.~Rev.}%          % Space Science Reviews
%\def\zap{ZAp}%          % Zeitschrift fuer Astrophysik
%\def\nat{Nature}%          % Nature
%\def\iaucirc{IAU~Circ.}%          % IAU Cirulars
%\def\aplett{Astrophys.~Lett.}%          % Astrophysics Letters
%\def\apspr{Astrophys.~Space~Phys.~Res.}%          % Astrophysics Space Physics Research
%\def\bain{Bull.~Astron.~Inst.~Netherlands}%          % Bulletin Astronomical Institute of the Netherlands
%\def\fcp{Fund.~Cosmic~Phys.}%          % Fundamental Cosmic Physics
%\def\gca{Geochim.~Cosmochim.~Acta}%          % Geochimica Cosmochimica Acta
%\def\grl{Geophys.~Res.~Lett.}%          % Geophysics Research Letters
%\def\jcp{J.~Chem.~Phys.}%          % Journal of Chemical Physics
\def\jgr{J.~Geophys.~Res.}%          % Journal of Geophysics Research
%\def\jqsrt{J.~Quant.~Spec.~Radiat.~Transf.}%          % Journal of Quantitiative Spectroscopy and Radiative Trasfer
%\def\memsai{Mem.~Soc.~Astron.~Italiana}%          % Mem. Societa Astronomica Italiana
%\def\nphysa{Nucl.~Phys.~A}%          % Nuclear Physics A
%\def\physrep{Phys.~Rep.}%          % Physics Reports
%\def\physscr{Phys.~Scr}%          % Physica Scripta
\def\planss{Planet.~Space~Sci.}%          % Planetary Space Science
%\def\procspie{Proc.~SPIE}%          % Proceedings of the SPIE
%\let\astap=\aap
%\let\apjlett=\apjl
%\let\apjsupp=\apjs
%\let\applopt=\ao
%

\usepackage{tikz}
\newcommand\tstrut{\rule{0pt}{2.5ex}}

%%%%%%%%%%%%%%%%%%%%%%%
%% Elsevier bibliography styles
%%%%%%%%%%%%%%%%%%%%%%%
%% To change the style, put a % in front of the second line of the current style and
%% remove the % from the second line of the style you would like to use.
%%%%%%%%%%%%%%%%%%%%%%%

%% Numbered
%\bibliographystyle{model1-num-names}

%% Numbered without titles
%\bibliographystyle{model1a-num-names}

%% Harvard
\bibliographystyle{model2-names.bst}\biboptions{authoryear}

%% Vancouver numbered
%\usepackage{numcompress}\bibliographystyle{model3-num-names}

%% Vancouver name/year
%\usepackage{numcompress}\bibliographystyle{model4-names}\biboptions{authoryear}

%% APA style
%\bibliographystyle{model5-names}\biboptions{authoryear}

%% AMA style
%\usepackage{numcompress}\bibliographystyle{model6-num-names}

%% `Elsevier LaTeX' style
%\bibliographystyle{elsarticle-num}
%%%%%%%%%%%%%%%%%%%%%%%
% EGGL
\usepackage[super]{nth}
\usepackage{fancyhdr}

\cfoot{\small{\textcopyright $\;$ 2019. California Institute of Technology. Government sponsorship acknowledged.}}
\pagestyle{fancy}
%%%%%%%%%%%%%%%%%%%%%%%
\begin{document}

\begin{frontmatter}

\title{Star catalog position and proper motion corrections in
asteroid astrometry II: The Gaia era \\ \textnormal{\normalsize 小行星天体测量中的星表位置和自行修正 II：Gaia 时代}}

%% Group authors per affiliation:
%\author{Elsevier\fnref{myfootnote}}
%\address{Radarweg 29, Amsterdam}
%\fntext[myfootnote]{Since 1880.}

%% or include affiliations in footnotes:
\author[JPL]{Siegfried Eggl\corref{mycorrespondingauthor}\fnref{uw}}
\cortext[mycorrespondingauthor]{Corresponding author}
\ead{eggl@uw.edu}

\author[JPL]{Davide Farnocchia}
\author[JPL]{Alan B. Chamberlin}
\author[JPL]{Steven R. Chesley}


\address[JPL]{Jet Propulsion Laboratory, California Institute of Technology, 4800 Oak Grove Drive, 91101 Pasadena, CA, USA}
\fntext[uw]{Present address: LSST/DIRAC Institute, Department of Astronomy, University of Washington, 15th Ave NE, Seattle, WA 98195, USA}

%\author[mymainaddress,mysecondaryaddress]{Elsevier Inc}
%\ead[url]{www.elsevier.com}

%\author[mysecondaryaddress]{Global Customer Service\corref{mycorrespondingauthor}}
%\cortext[mycorrespondingauthor]{Corresponding author}
%\ead{support@elsevier.com}

%\address[mymainaddress]{1600 John F Kennedy Boulevard, Philadelphia}
%\address[mysecondaryaddress]{360 Park Avenue South, New York}

\begin{abstract}
Astrometric positions of moving objects in the Solar System have been measured using a variety of star catalogs in the past.
Previous work has shown that systematic errors in star catalogs can affect the accuracy of astrometric observations. That, in turn, can influence the resulting orbit fits for minor planets. In order to quantify these systematic errors, we compare the positions and proper motion of stellar sources in the most utilized star catalogs to the second release of the Gaia star catalog. The accuracy of Gaia astrometry allows us to unambiguously identify local biases and derive a scheme that can be used to correct past astrometric observations of solar system objects.

过去，太阳系移动天体的天体测量位置是使用多种不同的星表测量出来的。之前的工作表明，星表中的系统误差会影响天体测量观测的精度。这反过来又会影响小行星的轨道拟合结果。为了量化这些系统误差，我们将最常使用的星表中的恒星位置和自行与 Gaia 星表的第二次发布数据进行了比较。Gaia 天体测量数据的精度使我们能够明确地识别局部偏差，并推导出一种方案，用于修正太阳系天体过去的天体测量观测数据。

Here we provide a substantially improved debiasing scheme for 26 astrometric catalogs that were extensively used in minor planet astrometry.
Revised corrections near the galactic center eliminate artifacts that could be traced back to reference catalogs used in previous debiasing schemes.

在这里，我们为 26 个在小行星天体测量中广泛使用的星表提供了一个大幅改进的去偏方案。修订后的银河系中心附近的修正消除了以前去偏方案中因参考星表而产生的伪影。

Median differences in stellar positions between catalogs now tend to be on the order of several tens of milliarcseconds (mas) but can be as large as 175 mas. Median stellar proper motion corrections scatter around 0.3 mas/yr and range from 1--4 mas/yr for star catalogs with and without proper motion, respectively.

星表之间恒星位置差异的中位数现在通常在几十毫角秒（mas）的数量级，但在某些情况下可高达 175 mas。恒星自行修正的中位数分布在 0.3 mas/yr 左右，对于包含和不包含自行的星表，其修正范围分别为 1-4 mas/yr。

The tables in this work are meant to be applied to existing optical observations. They are not intended to correct new astrometric measurments as those should make use of the Gaia astrometric catalog.
Since previous debiasing schemes already reduced systematics in past observations to a large extent, corrections beyond the current work may not be needed in the foreseeable future.

本工作中的表格旨在应用于现有的光学观测数据。它们不打算用于修正新的天体测量结果，因为新的测量应直接使用 Gaia 天体测量星表。由于以前的去偏方案已经在很大程度上减少了过去观测中的系统误差，在可预见的未来，可能不需要进行超出本工作范围的修正。

.   %Zonal differences are .
% Correcting past observations based on Gaia DR 2 will only slightly improve our ability to predict the future whereabouts of Solar System objects.
\end{abstract}

\begin{keyword}
Asteroids --- Asteroids, dynamics --- Orbit determination --- Astrometry --- Near-Earth objects
\end{keyword}

\end{frontmatter}
%%%
%\linenumbers
%%%
\section{Introduction}
\textbf{【译】引言}

Over more than two centuries astronomers have collected approximately 200 million astrometric observations of minor planets in the Solar System\footnote{Minor Planet Center (MPC), http://www.minorplanetcenter.net/}.
This dataset %trove 
of astrometric measurements is essential to estimate the trajectories of Solar System Objects (SSOs) such as near-Earth asteroids \citep{milani10}.

两个多世纪以来，天文学家收集了大约 2 亿次太阳系小行星的天体测量观测数据。这一庞大的天体测量数据集对于估算太阳系天体（SSOs），如近地小行星的轨迹至关重要（Milani 和 Gronchi, 2010）。

The vast majority of astrometry used for orbit determination of SSOs is optical, i.e. pairs of Right Ascension (RA) and Declination (DEC) measurements at given epochs. In contrast, only a tiny fraction of the SSO population has been observed by radar\footnote{About one in a thousand asteroids, see NASA Jet Propulsion Laboratory (JPL), https://echo.jpl.nasa.gov/asteroids/index.html/}. Improving the quality of optical astrometry is a worthwhile endeavor, especially if that means extending the arc of observations for objects of interest.  

用于 SSO 轨道确定的绝大多数天体测量数据都是光学的，即给定历元的赤经（RA）和赤纬（DEC）对。相比之下，只有极小部分的 SSO 群体被雷达观测过。提高光学天体测量的质量是一项值得努力的工作，特别是如果这意味着能够延长感兴趣目标的观测弧段。
%
Over the years a variety of star catalogs has been used to measure astrometric positions of asteroids to a common frame of reference such as the ICRF \citep{icrf2} (Table \ref{tab:catalogs}).
With improving quality of astrometric measurements it became evident that some observations of asteroids showed systematic errors that deviated from the expected scatter around the nominal trajectory \citep{carpino03}.  
Some of those systematic errors could be traced back to the use of specific astrometric catalogs in the observation reduction process. 

多年来，人们使用了各种星表来测量小行星的天体测量位置，使其统一到一个共同的参考框架下，例如 ICRF（Ma 等人，2009）（表 1）。随着天体测量精度的提高，人们发现小行星的一些观测结果显示出系统误差，这些误差偏离了标称轨迹周围的预期离散度（Carpino 等人，2003）。其中一些系统误差可以追溯到在观测归算过程中使用了特定的天体测量星表。

To assess these systematic errors one can compare stellar positions and proper motion between various star catalogs. 
\cite{chesley10} used the 2MASS \citep{2mass} catalog as a reference to identify local systematics in positions of stars in the USNO catalogs \citep{usnoa2,usnoa1,usnob1}. Correcting for regional biases in star positions led to better ephemeris predictions and improved the error statistics of astrometric observations of asteroids.

为了评估这些系统误差，我们可以比较不同星表之间的恒星位置和自行。Chesley 等人 (2010) 使用 2MASS (Skrutskie 等人，2006) 星表作为参考，识别了 USNO 星表 (Monet, 1998; Lasker 等人, 1996a; Monet 等人, 2003) 中恒星位置的局部系统误差。修正恒星位置的区域性偏差使得星历预报更加准确，并改善了小行星天体测量观测的误差统计。

Selecting the most accurate stars in the PPM XL catalog based on 2MASS astrometry, \cite{farnocchia15} improved the scheme by \cite{chesley10} and debiased practically all of the prevalent star catalogs used for asteroid astrometry. \cite{farnocchia15} included proper motion in the debiasing process which led to widespread improvements in the post-fit residuals of asteroids, most notably (99942) Apophis, (101955) Bennu and (6489) Golevka. 

Farnocchia 等人 (2015) 基于 2MASS 天体测量数据选择了 PPM XL 星表中最精确的恒星，改进了 Chesley 等人 (2010) 的方案，并对几乎所有用于小行星天体测量的流行星表进行了去偏。Farnocchia 等人 (2015) 在去偏过程中加入了自行修正，这导致小行星的拟合后残差得到了广泛的改善，最显著的是 (99942) Apophis、(101955) Bennu 和 (6489) Golevka。 

%
The present work follows a similar approach with the advantage of having Gaia data, described in section \ref{sec:gaiadr2}, as a reference. 
Following a brief outline of the debiasing process (section \ref{sec:debias}) we calculate position and proper motion corrections for stellar catalogs that have been used extensively in minor planet astrometry (sections \ref{sec:cats} - \ref{sec:interp}). 
In section \ref{sec:eph} we then assess the quality of ephemeris predictions based on the new debiasing scheme.
A discussion of our results concludes this article.

目前的工作采用了类似的方法，但优势在于拥有第 2 节中描述的 Gaia 数据作为参考。在简要概述去偏过程（第 3 节）之后，我们计算了广泛用于小行星天体测量的星表的位置和自行修正（第 4-6 节）。随后在第 7 节中，我们评估了基于新去偏方案的星历预报质量。最后我们在文章结尾讨论了我们的结果。
%
\begin{table}\tiny
\begin{center}
\begin{tabular}{llccccc}
\hline
Catalog & \multicolumn{2}{c}{MPC designation}  & \multicolumn{2}{c}{Asteroid observations} & VizieR & Reference\\
             & ADES  & flag & count &  \%   &  CDS ID &\\
\hline
USNO A2.0   & USNOA2  & c  & 40,964,725 & 21.929 & I/252   &\citet{usnoa2}\\
2MASS       & 2MASS  & L  & 36,656,505 & 19.623 &  II/246   & \citet{2mass}\\
UCAC 2      &  UCAC2 & r  & 30,808,742 & 16.493 & I/289    &\citet{ucac2}\\
UCAC 4      & UCAC4  & q  & 25,280,408 & 13.533 & I/322A       &\citet{ucac4}\\
USNO B1.0   & USNOB1  & o  & 16,827,916 & 9.008 & I/284       &\citet{usnob1}\\
SST RC4     &  SSTRC4 & R  & 14,143,319 & 7.571 & N/A       & \citet{sstrc5} \\
Gaia DR 1    & Gaia1  & U  & 7,878,286 & 4.217 & I/337         &\citet{gaiadr1}\\
UCAC 3      &  UCAC3 & u  & 3,460,491 & 1.852 & I/315         &\citet{ucac3}\\
USNO A1.0   & USNOA1 & a  & 2,196,215 & 1.175 & N/A         &\citet{usnoa1}\\
Unspecified     & N/A & N/A & 1,995,344 & 1.068 & N/A       &N/A\\
USNO SA2.0  & USNOSA2 & d  & 1,757,923 & 0.941 & N/A        &\citet{usnoa2}\\
SDSS DR 7    & SDSS7  & N  & 985,016 & 0.527 & II/294          &\citet{sdss7}\\
GSC 1.1     &  GSC1.1 & i  & 645,795 & 0.345 & I/220          &\citet{gsc1.1}\\
UCAC 1      & UCAC1  & e  & 516,449 & 0.276 &   I/268        &\citet{ucac1}\\
GSC ACT     & GSCACT  & m  & 474,122 & 0.253 & I/255         &\citet{gsc_act}\\
CMC 14      & CMC14  & w  & 426,781 & 0.228 & I/304          &\citet{cmc14}\\
Tycho 2     & Tyc2  & g  & 414,400 & 0.221 &  I/259          &\citet{tycho2}\\
USNO SA1.0  & USNOSA1  & b  & 349,099 & 0.186 & N/A        &\citet{usnoa1}\\
PPM XL       & PPM XL  & t  & 302,432 & 0.162 & I/317         &\citet{ppmxl}\\
GSC (unspec.)& GSC & z & 292,814 & 0.156    & N/A          &N/A\\
ACT         & ACT  & l & 112,431 & 0.060 & I/246             &\citet{act}\\
Other catalogs & N/A &N/A& 97,256 & $0.052$ & N/A & N/A\\
NOMAD       & NOMAD  & v  & 82,997 & 0.044 & N/A           &\citet{nomad}\\
PPM         & PPM  & p & 47,227 & 0.025 & I/146,I/193              &\citet{ppm}\\
%Yale        &  Yale & K & 37,672 & 0.020 & I/238A             &\citet{yale}\\
URAT 1      & URAT1  & S & 37,105 & 0.019 & I/329            &\citet{urat1}\\
%SAO         & SAO  & C & 28,964 & 0.015 & I/131A              &\citet{sao}\\
GSC 1.2     & GSC1.2& j  & 17,269 & 0.009 & I/254            &\citet{gsc1.2}\\
%MPOSC3      & MPOSC3  & P & 14,757 & 0.008 &  N/A          &\citet{mposc3}\\
CMC 15      & CMC15& Q  & 11,869 & 0.006 &  I/327            &\citet{cmc15}\\
%AC          & AC  & A  & 7,071 & 0.004 & I/96                &\citet{ac}\\
SDSS DR 8    & SDSS8  & n &  6,157 & 0.003 & II/306            &\citet{sdss8}\\
%AGK 3        & AGK3  & D &  4,659 & 0.002 &  I/61B            &\citet{agk3}\\
UCAC 5   & UCAC5 & W & 1,618 & $<$0.001 & I/340 & \citet{ucac5} \\
%Hipparcos 2& Hip2 & x & 118 & $<$0.001 & I/311 & \citet{hip2}  \\
%GSC-2.2 & k & & 745 & 0.00 & \citet{gsc2.2}\\
%GSC-2.3 & M & & 86 & 0.00 & \citet{gsc2.3}\\

\hline
\end{tabular}
\end{center}
\caption{Star catalogs considered in this work. This table accounts for all the astrometry
  available up to March 23, 2018.}
\label{tab:catalogs}
\end{table}
%%%
\section{The Gaia catalogs} 
\textbf{【译】Gaia 星表}
\label{sec:gaiadr2}
Since 2013 the Gaia mission has been collecting astrometric information on roughly 1.7 billion sources \citep{gaiadr2}. A first data release based on a partially completed survey featured around a billion sources largely without proper motion \citep{gaiadr1}.
With the second Gaia data release (Gaia DR 2) published on April 25, 2018 this information has become publicly accessible through the Gaia Archive\footnote{Gaia Archive, https://gea.esac.esa.int/archive/}. 

自 2013 年以来，Gaia 任务一直在收集大约 17 亿个源的天体测量信息 (Gaia Collaboration 等人, 2018)。基于部分完成的巡天的第一次数据发布包含了大约 10 亿个源，但大部分没有自行数据 (Lindegren 等人, 2016)。随着 2018 年 4 月 25 日发布的第二次 Gaia 数据（Gaia DR 2），这些信息已通过 Gaia 档案公开发布。

%The high uniformity and excellent quality of Gaia DR 2 parallaxes and proper motions could be achieved by relying exclusively on data from the Gaia satellite. 
About 1.3 billion sources between 3$<$Gmag$<$21 come with a full five-parameter astrometric solution, i.e. right ascension and declination positions on the sky $(\alpha,\delta)$, parallaxes, and proper motion. Uncertainties in the parallaxes range between 0.04 milliarcseconds (mas) for sources with Gmag$<$15 and 0.7 mas at Gmag$=$20. The corresponding uncertainties in the proper motion components start at 0.06 mas/yr for Gmag$<$15 and reach 1.2 mas/yr for Gmag$=$20. The Gaia astrometric catalog is, thus, a reference of unprecedented accuracy well suited to identify systematic errors in other star catalogs.

大约 13 亿个亮度在 3$<$Gmag$<$21 之间的源拥有完整的五参数天体测量解，即天空中的赤经和赤纬位置 ($\alpha$, $\delta$)、视差和自行。对于 Gmag$<$15 的源，视差的不确定度在 0.04 毫角秒（mas）之间，而对于 Gmag$=$20 的源，不确定度为 0.7 mas。相应的自行分量不确定度对于 Gmag$<$15 的源从 0.06 mas/yr 开始，对于 Gmag$=$20 的源达到 1.2 mas/yr。因此，Gaia 天体测量星表是一个具有前所未有精度的参考标准，非常适合用于识别其他星表中的系统误差。
%%%
\section{Star position and proper motion corrections} 
\textbf{【译】星表位置和自行修正}
\label{sec:debias}
Similar to \citet{chesley10} and \citet{farnocchia15} we compare the various star catalogs that were used for minor planet astrometry in the past to a reference catalog, in this case Gaia DR 2. 
% Potential catalog systematics are known to vary across the sky \citep{farnocchia15}. (XXX I would remove this sentence, also this was known to Chesley et al. 2010 too XXX)
By dividing the celestial sphere into 49,152
equal-area tiles ($\sim$0.8 square degrees each) using a HEALPix tesselation
\citep{healpix} and comparing stellar positions and proper motion within a given tile we can quantify local systematic differences.

与 Chesley 等人 (2010) 和 Farnocchia 等人 (2015) 类似，我们将过去用于小行星天体测量的各种星表与参考星表（在本例中为 Gaia DR 2）进行比较。通过使用 HEALPix 镶嵌 (Gorski 等人, 2005) 将天球划分为 49,152 个等面积图块（每个约 0.8 平方度），并比较给定图块内的恒星位置和自行，我们可以量化局部的系统性差异。

To this end we extract stellar positions and proper motion at epoch J2000 from the reference and the test catalog.
Given the variability in the number of stars per tile in each catalog, identification of common stars requires some attention. We first search for spatial correlation between astrometric sources within  1'' of the reference position. 
If candidates are found we apply the identification and outlier rejection criterion described in \citet{farnocchia15} to
ensure spurious matches are excluded from the sample.

为此，我们提取了参考星表和测试星表在 J2000 历元的恒星位置和自行。考虑到每个星表在每个图块中的恒星数量具有可变性，识别共同的恒星需要一些关注。我们首先搜索参考位置 1 角秒范围内的天体测量源的空间相关性。如果找到候选者，我们应用 Farnocchia 等人 (2015) 中描述的识别和异常值剔除标准，以确保从样本中排除虚假匹配。

%
The median of the differences in the astrometric quantities of each identified star per tile yields the local corrections for position $(\Delta \alpha_{J2000}, \Delta \delta_{J2000})$, and proper motion $(\Delta \mu_{\alpha}, \Delta \mu_{\delta}$) at epoch J2000.
%
Astrometric observations of minor planets can then be corrected for catalog systematics by subtracting
 \begin{eqnarray*}
\Delta\alpha & = & \Delta \alpha_{J2000} + \Delta \mu_{\alpha} \cdot \Delta t,\\
\Delta\delta & = & \Delta \delta_{J2000} + \Delta \mu_{\delta} \cdot \Delta t,\\
\end{eqnarray*}
with $\Delta t=t - 2451545.0$, the difference between the epoch of observation and J2000 in JD.
All differential quantities $\Delta{\alpha}$, $\Delta{\alpha}_{J2000}$, and $\Delta \mu_{\alpha}$ include the spherical metric factor $\cos\delta$.

每个图块中已识别恒星的天体测量量差异的中位数，产生了 J2000 历元的局部位置修正 $(\Delta \alpha_{J2000}, \Delta \delta_{J2000})$ 和自行修正 $(\Delta \mu_{\alpha}, \Delta \mu_{\delta})$。然后，可以通过减去以下各项来修正小行星天体测量观测中的星表系统误差：
(公式略)
其中 $\Delta t=t - 2451545.0$，即观测历元与 J2000 之间的儒略日（JD）差值。所有的差分量 $\Delta{\alpha}$、$\Delta{\alpha}_{J2000}$ 和 $\Delta \mu_{\alpha}$ 都包含了球面度规因子 $\cos\delta$。
%
%In the following sections we briefly introduce the set of astrometric catalogs that has not been discussed in \citet{farnocchia15}.
%%%
\section{Added astrometric catalogs} 
\textbf{【译】新增的天体测量星表}
\label{sec:cats}
In addition to the 20 catalogs discussed in \citet{farnocchia15} and Gaia DR 1 we added the following five catalogs to be compared to Gaia DR 2:
% CMC 15, %Hipparcos 2, 
% SDSS DR 8, SST RC 5, UCAC 5, URAT 1  and PPM XL itself to the comparison with Gaia DR 2.

除了 Farnocchia 等人 (2015) 中讨论的 20 个星表和 Gaia DR 1 之外，我们还添加了以下五个星表与 Gaia DR 2 进行比较：
\begin{itemize}
\item CMC 15 is the \nth{15} installment of the Carlsberg Meridian Catalogue generated through the Carlsberg Meridian Telescope, formerly the Carlsberg Automatic Meridian Circle \citep{cmc15}. CMC 15 has been compiled between March 1999 and March 2011, encompasses around 134,000,000 entries, and claims a positional accuracy between 35 to 100 mas for a declination range from -40 and +50 degrees.
% Hipparcos 2 was a rereduction of the Hipparcos catalog published in 1997 providing a substantial improvement - about a factor of two - on the already excellent astrometry (target accuracy 2 mas) delivered by ESA's Hipparcos mission \citep{hip2}. One of the limitations of the Hippachos 2 catalog, and the main reason why it has not been more popular in minor planet astrometry, is its relatively small size. Hipparcos 2 contains merely 117,955 sources.    

\item The \nth{8} data release of the Sloan Digital Sky Survey (SDSS) added roughly 4 times the number of sources published in data release 7, namely about 325 million \citep{sdss8}. The fact that astrometric data together with proper motion is available has made SDSS DR 7 \citep{sdss7} attractive for minor planet astrometry. Almost 1 million observations of minor planets were reduced with SDSS DR 7 (Table \ref{tab:catalogs}). SDSS DR 8, although offering a significant increase in source density, has not been as popular, which is partly due to known issues with the astrometry \citep{sdss8_astrometry}.         

\item The Space Surveillance Telescope (SST) has been relying on a proprietary catalog based on selected sources from a variety of other astrometric and photometric catalogs for the reduction of its 14+ million observations of asteroids in the Solar System. Although not in the public domain, the \nth{5} version of the SST catalog (SST-RC5) has been kindly made available for this study \citep{sstrc5}. Most of the astrometric observations considered in this work have been submitted under SST-RC4. Astrometric standard stars in SST-RC5 are largely the same as in the \nth{4} installment of the SST catalog, however \citep{sstrc5}. Hence, we use the same corrections to debias astrometric obsevations submitted under SST-RC4 and SST-RC5. 

\item UCAC 5, is the \nth{5} installment of the US Naval Observatory CCD Astrograph Catalog (UCAC) \citep{ucac5}. By combining UCAC 5 with Gaia DR1 data new proper motions on the Gaia coordinate system for over 107 million stars were obtained with typical accuracies of 1--2 mas/yr for R=11--15 mag, and about 5 mas/yr through R=16 mag.

\item URAT, a follow-up project to the UCAC series featured and increased limiting magnitude that allowed for a roughly 4-fold increase in the average number of stars per square degree as compared to UCAC. Additionally, stars as bright as about \nth{3} magnitude were added to the catalog. The URAT 1 catalog \citep{urat1} has its mean epoch between 2012.3 and 2014.6 supporting a magnitude range from 3 to 18.5 in R-band with a positional precision of 5 to 40 mas. It covers most of the northern hemisphere and some areas down to -24.8 deg in declination.
\end{itemize}
%%%
\section{Catalog Comparison Results}
\textbf{【译】星表比较结果}
\label{sec:catresults}

With more than 40 million entries in the Minor Planet Center database the USNO A2.0 catalog has been one of the most popular choices for reducing astrometric observations of minor planets.
Figures \ref{fig:usnoa2} and \ref{fig:usnoa2_hist} show the proposed corrections derived from a comparison with Gaia DR 2.
The regional biases in stellar positions first found by \citet{chesley10} and later confirmed by \citet{farnocchia15} are clearly visible in the top panels of Figure \ref{fig:usnoa2}.

在小行星中心数据库中拥有超过 4000 万个条目，USNO A2.0 星表一直是归算小行星天体测量观测结果最流行的选择之一。图 1 和图 2 显示了通过与 Gaia DR 2 比较得出的建议修正。Chesley 等人 (2010) 首先发现并后来由 Farnocchia 等人 (2015) 确认的恒星位置区域偏差，在图 1 的顶部面板中清晰可见。

%
Figure \ref{fig:usnoa2_hist} shows the distribution of the corrections collected over the entire sky. The visual impression from Figure \ref{fig:usnoa2}, namely that there is an overall bias towards higher declination values in stellar positions, is confirmed in the distribution of $\Delta \delta_{J2000}$. The lack of proper motion as well as a median bias of 162 mas in declination with a mode closer to a quarter arcsecond all suggest that astrometric observations reduced with USNO A2.0 could benefit from a correction. 

图 2 显示了在整个天球上收集的修正分布。从图 1 中得到的视觉印象，即恒星位置普遍偏向较高的赤纬值，在 $\Delta \delta_{J2000}$ 的分布中得到了确认。缺乏自行以及赤纬方向 162 mas 的中值偏差（其众值接近四分之一个角秒）都表明，使用 USNO A2.0 归算的天体测量观测可以从修正中受益。
%
Table~\ref{tab:cat_stats} reports the size of the corrections in terms of the all sky median (MED) and median absolute deviation (MAD) for each studied catalog. 

表 2 报告了每个研究星表的全球中值（MED）和中值绝对偏差（MAD）的修正量。

%
Of particular interest is the comparison between PPM XL and Gaia DR 2, since a subset of the former was used as a reference for stellar positions and proper motion in \citet{farnocchia15}.
The corresponding graphs are presented in Figure \ref{fig:PPMXL}.
The range of position corrections is only about one sixth of that in USNO A2.0 and the subset of 2MASS stars used for astrometry also has a relatively small positional bias as can be gathered from Table \ref{tab:cat_stats}.

特别令人感兴趣的是 PPM XL 与 Gaia DR 2 之间的比较，因为前者的一个子集曾在 Farnocchia 等人 (2015) 中被用作恒星位置和自行的参考。相应的图表见图 3。位置修正的范围只有 USNO A2.0 的约六分之一，从表 2 中可以看出，用于天体测量的 2MASS 恒星子集也具有相对较小的位置偏差。

%
A comparison between Gaia DR 2 and UCAC 4 (Figure \ref{fig:ucac4}) reveals that
the relatively large corrections seen by \citet[][Figures 9 to 12]{farnocchia15} near the galactic center ($\alpha\approx$ 270 deg, $\delta\approx$-30 deg) appear to actually be associated to the PPM XL reference, rather than the UCAC catalogs. This is indicated by the signature seen near the galactic center in all four panels of Figure \ref{fig:PPMXL} that is now substantially absent from all four panels of Figure \ref{fig:ucac4}.

Gaia DR 2 与 UCAC 4 的比较（图 5）表明，Farnocchia 等人 (2015, 图 9 到 12) 在银河系中心附近 ($\alpha\approx$ 270 度, $\delta\approx$-30 度) 看到的相对较大的修正，实际上似乎与 PPM XL 参考星表有关，而不是 UCAC 星表。图 3 的所有四个面板中银河系中心附近看到的特征，在图 5 的所有四个面板中现已基本消失。

% There is, however, a significant non-uniformity in the corrections close to the Galactic center ($\alpha\approx$ 270 deg, $\delta\approx$-30 deg) both in position and proper motion, which leads to a drop in the fraction of stars that could be matched between Gaia DR 2 and PPM XL. Despite the careful selection of the subset of stars that was used as a reference in \citet{farnocchia15}---2MASS has very little bias---the non-uniformity in PPM XL affected the previous debiasing scheme. This signature is most evident in the corrections for the UCAC catalog series as presented in \citet[][Figures 9 to 12]{farnocchia15}.
% %
% A comparison between Gaia DR 2 and UCAC 4 (Figure \ref{fig:ucac4}) does not show any regional bias near the Galactic center in UCAC 4, although proper motion bias is generally higher in the vicinity of the Milky Way.
%

% Comparing the distributions of PPM XL and UCAC 4 catalog corrections in Figures \ref{fig:PPM XL_hist} and \ref{fig:ucac4_hist} one can see that the latter shows less bias. Similar conclusions hold for UCAC 2 as can be seen in Table \ref{tab:cat_stats}. Both UCAC 2 and UCAC 4 were used extensively in minor planet astrometry totalling around 56 million observations.

With the exception of the artificial structures close to the Galactic center introduced through PPM XL, our results for UCAC, USNO and CMC catalogs are very similar to those found in \citet{farnocchia15}. 
The corresponding Figures are provided in the Appendix.

除了通过 PPM XL 引入的银河系中心附近的人为结构外，我们对 UCAC、USNO 和 CMC 星表的结果与 Farnocchia 等人 (2015) 发现的结果非常相似。相应的图表见附录。

%
Among the newly added catalogs UCAC 5 continues the tradition of the UCAC family with an increase in the overall quality of the catalog. Local corrections are mostly due to differences between stellar positions in Gaia DR 2 and Gaia DR 1. The latter were ingested into UCAC 5. URAT 1 has a higher source density compared to UCAC 4, but it also exhibits larger deviations from Gaia DR 2 in both stellar positions and proper motion. 

在新增的星表中，UCAC 5 延续了 UCAC 家族的传统，星表的整体质量有所提高。局部修正主要是由于 Gaia DR 2 和 Gaia DR 1 之间的恒星位置差异，后者被吸纳进了 UCAC 5 中。与 UCAC 4 相比，URAT 1 具有更高的源密度，但在恒星位置和自行方面偏离 Gaia DR 2 的程度也更大。

%
We confirm the known issues with astrometric accuracy around the north celestial pole in SDSS DR 8 \citep{sdss8_astrometry}, which were addressed in the following data releases.  

我们确认了 SDSS DR 8 在北天极附近天体测量精度的已知问题 (Aihara 等人, 2011)，这些问题在随后的数据发布中得到了解决。

%
Apart from minor systematics stemming from survey incompleteness, Gaia DR 1 had excellent stellar positions at the epoch of its publication \citep{gaiadr1}. Since Gaia DR 1 largely lacks proper motion, the differences to Gaia DR 2 become significant at the epoch of comparison (J2000), however. This is reflected in Table \ref{tab:cat_stats}. The quality of stellar positions, even if measured through the Gaia satellite, degrades relatively quickly with time if no proper motion corrections are applied. Minor planet astrometrists are, thus, encouraged to switch to Gaia DR 2 as soon as possible. 

除了巡天不完整引起的小规模系统误差外，Gaia DR 1 在其发布历元时具有极好的恒星位置。然而，由于 Gaia DR 1 很大程度上缺乏自行，因此在比较历元（J2000）时，与 Gaia DR 2 的差异变得显著（见表 2）。即使是由 Gaia 卫星测量的恒星位置质量，如果不应用自行修正，也会随着时间推移而下降得相对较快。因此，鼓励小行星天体测量学家尽快转向使用 Gaia DR 2。

%
We debias all observations where corrections are available with the exception of those reduced with ACT, Gaia DR 1, Tycho-2 and UCAC-5, since the latter show no large scale structures in local position and proper motion systematics.

我们对所有有修正数据的观测进行去偏，但使用 ACT、Gaia DR 1、Tycho-2 和 UCAC-5 归算的观测除外，因为后者在局部位置和自行系统误差方面没有表现出大尺度结构。
%
\begin{figure}
\begin{tabular}{ll}
\includegraphics[width=0.5\linewidth]{figs/usno_a2_0.png} &
\includegraphics[width=0.5\linewidth]{figs/usno_a2_1.png} \\
\includegraphics[width=0.5\linewidth]{figs/usno_a2_2.png} &
\includegraphics[width=0.5\linewidth]{figs/usno_a2_3.png}
\end{tabular}
\caption{USNO A2.0 catalog systematics with respect to Gaia DR 2. The top panels display local differences in right ascension and declination, the bottom panels show the systematics caused by the missing proper motion.   \label{fig:usnoa2}}
\end{figure}
%
\begin{figure}
\includegraphics[width=\linewidth]{figs/usno_a2_hist.png} 
\caption{Summary statistics of USNO A2.0 catalog systematics with respect to Gaia DR 2. The top panel displays the distribution of local differences in right ascension and declination. The signature of the missing proper motion is shown in the bottom panel. The legend contains the median (MED) and median absolute deviation (MAD) for the respective quantities.\label{fig:usnoa2_hist}}
\end{figure}
%
\begin{figure}
\begin{tabular}{ll}
\includegraphics[width=0.5\linewidth]{figs/ppmxl_0.png} &
\includegraphics[width=0.5\linewidth]{figs/ppmxl_1.png} \\
\includegraphics[width=0.5\linewidth]{figs/ppmxl_2.png} &
\includegraphics[width=0.5\linewidth]{figs/ppmxl_3.png}\\
% \includegraphics[width=0.5\linewidth]{figs/ppmxl_0_nstar.png} &
% \includegraphics[width=0.5\linewidth]{figs/ppmxl_fmatch.png} 
\end{tabular}
\caption{Same as Figure \ref{fig:usnoa2} but for the
PPM XL catalog. 
% In addition, the number of stars per HEALPix tile in PPM XL (bottom left) and fraction of matching stars between PPM XL and Gaia DR 2 (bottom right) are shown.
%Note the drop in the fraction of matching stars near the galactic center.
Note the larger corrections near the galactic center ($\alpha\approx$ 270 deg, $\delta\approx$-30 deg).
\label{fig:PPMXL}}
\end{figure}
%
\begin{figure}
\includegraphics[width=\linewidth]{figs/ppmxl_hist.png} 
\caption{Same as Figure \ref{fig:usnoa2_hist} for the PPM XL catalog.  \label{fig:ppmxl_hist}}
\end{figure}
%
\begin{figure}
\begin{tabular}{ll}
\includegraphics[width=0.5\linewidth]{figs/ucac4_0.png} &
\includegraphics[width=0.5\linewidth]{figs/ucac4_1.png} \\
\includegraphics[width=0.5\linewidth]{figs/ucac4_2.png} &
\includegraphics[width=0.5\linewidth]{figs/ucac4_3.png}\\
% \includegraphics[width=0.5\linewidth]{figs/ucac4_0_nstar.png} &
% \includegraphics[width=0.5\linewidth]{figs/ucac4_fmatch.png}
\end{tabular}
\caption{Same as Figure \ref{fig:PPMXL} but for the UCAC 4 catalog. Systematics are small in positions, but traces of stellar proper motion errors in the galactic plane are visible. \label{fig:ucac4}}
\end{figure}
%
\begin{figure}
\includegraphics[width=\linewidth]{figs/ucac4_hist.png} 
\caption{Same as Figure \ref{fig:usnoa2_hist} for the UCAC 4 catalog. \label{fig:ucac4_hist}}
\end{figure}
%
\begin{table}\tiny
\begin{center}
\begin{tabular}{|l|cc|cc|cc|cc|c|c|}
  \hline
Catalog & \multicolumn{2}{c}{$\Delta{\alpha}_{J2000}$ [mas]} & \multicolumn{2}{c}{$\Delta{\delta}_{J2000}$ [mas]} & \multicolumn{2}{c}{$\Delta\mu_{\alpha}$ [mas/yr]} &
\multicolumn{2}{c}{$\Delta\mu_{\delta}$ [mas/yr]} & Sky [\%] & proper\\
& MED & MAD & MED & MAD & MED & MAD &MED & MAD &coverage & motion\\ 
%& [arcsec] & [arcsec]  & [mas/yr] & [mas/yr] & \\
\hline
2MASS&  -2.4& 15.5& 11.6& 13.6& 1.51& 2.50& 3.76& 1.50& 100 & no \\ 
%AC &219.2& 304.0& 292.5& 246.1& 0.28& 1.93& 3.54& 1.55& 28  & no\\ 
ACT&  -0.8& 11.8& 0.2& 10.9& -0.05& 1.03& 0.01& 0.98& 100 & yes\\ 
CMC 14& 2.0& 16.7& -18.3& 13.8& 1.53& 2.45& 4.22& 0.98& 61& no \\ 
CMC 15&  1.5& 19.0& -19.8& 15.7& 1.52& 2.53& 4.09& 1.17&70& no \\ 
GSC 1.1& -45.5& 266.8& 17.1& 235.1& 1.07& 2.81& 3.87& 1.48&100& no \\ 
GSC 1.2& 23.6& 127.7& -16.2& 119.3& 1.06& 2.81& 3.87& 1.48& 100& no \\ 
GSC ACT& 29.4& 82.2& 44.8& 69.9& 1.05& 2.76& 3.84& 1.45& 100& no\\ 
Gaia DR 1& -22.5& 30.5& -49.2& 20.2& 1.48& 2.00& 3.27& 1.35& 100 & no \\ 
%Hipparcos 2& 0.1& 7.1& 0.1& 6.1& 0.12& 0.76& 0.01& 0.66& 100& yes\\ 
NOMAD& -19.8& 62.3& 154.8& 80.8& 1.19& 1.78& 2.66& 1.29& 100& yes\\ 
PPM& 12.6& 131.2& -68.4& 130.0& -0.21& 2.77& -0.48& 2.84& 100& yes\\ 
PPM XL& -19.5& 18.1& 36.0& 17.0& -0.30& 1.34& -0.23& 1.46& 100& yes\\ 
SDSS DR 7& -5.4& 16.4& -7.6& 16.2& 0.29& 0.35& 0.67& 0.43& 13& yes\\ 
SDSS DR 8& -4.7& 22.0& -5.2& 24.0& 0.30& 0.47& 0.75& 0.51& 16& yes\\ 
SST RC 4\&5&-15.2& 15.0& 25.0& 14.4& 0.27& 1.08& 0.10& 1.13& 100& yes\\ 
Tycho 2& -1.0& 10.1& -1.6& 9.8& -0.08& 0.82& -0.31& 0.82& 100& yes\\ 
UCAC 1& -4.0& 14.0& 12.5& 13.1& 2.70& 3.06& 2.94& 4.39& 39& yes\\ 
UCAC 2& -0.8& 10.3& 3.1& 9.2& -0.30& 1.31& -0.60& 1.27& 88& yes\\ 
UCAC 3& -2.3& 13.3& -1.6& 10.9& -0.75& 1.82& 0.84& 1.87& 100& yes\\ 
UCAC 4& 1.0& 11.0& -2.8& 9.6& 0.10& 0.84& 0.34& 0.94& 100& yes\\ 
UCAC 5& 1.6& 4.8& 1.0& 3.7& -0.11& 0.32& -0.06& 0.25& 100& yes\\ 
URAT 1& -23.0& 26.9& -43.7& 20.4& 0.42& 1.17& 0.70& 1.14& 61& yes\\ 
USNO A1.0&-17.3& 261.4& -22.9& 231.1& 1.39& 1.92& 3.21& 1.34& 100& no\\ 
USNO A2.0& 75.3& 111.9& 162.2& 132.4& 1.40& 1.91& 3.20& 1.34& 100& no\\ 
USNO B1.0& -25.4& 79.3& 175.2& 76.8& 1.26& 2.00& 3.02& 1.36& 100& yes\\ 
USNO SA1.0& -20.3& 260.8& -17.1& 230.8& 1.27& 1.96& 3.27& 1.34&100& no\\
USNO SA2.0& 72.7& 111.8& 164.1& 131.5& 1.35& 1.96& 3.25& 1.34& 100& no\\ 
\hline
\end{tabular}
\end{center}
\caption{The columns show the median (MED) and median absolute deviation (MAD) of corrections in
  position (right ascension and declination) and proper motion (right ascension and declination) over the entire sky as well as the fraction of sky coverage for each catalog. The last column states whether or not a catalog contains data on stellar proper motion.}
\label{tab:cat_stats}
\end{table}
%%%
\section{Interpolation of Debiasing Tables}
\textbf{【译】去偏表的插值}
\label{sec:interp}
Due to the discrete nature of calculating corrections to astrometric measurments on HEALPix tiles, transitions between neighboring corrections are not guaranteed to be smooth. 
As an example, Figure \ref{fig:smooth} shows that adjacent USNO A2.0 tiles can have significantly different bias values, with variations upwards of 0.6 arcsec.
% Figure \ref{fig:smooth} shows, for instance, that in USNO A1.0 tiles with positive and negative corrections on the order of an arcsecond can have common borders. 
Corrections for minor planet observations that cross such a border would see an artificial jump in the corresponding coordinate. A finer HEALPix tessellation could be used to mitigate such a problem. However, a finer tessellation would mean that fewer stars per tile are available to estimate the magnitude of the correction. 
In order to circumvent reducing the sample size of stars per tile, we use radial basis functions \citep[RBFs,][]{schaback1995} to interpolate data from neighboring tiles onto a finer HEALPix tessellation. This reduces artificial discontinuities along tile borders while at the same time retaining a sufficient number of stars in each of the stencil tiles.   

由于在 HEALPix 图块上计算天体测量差值的离散性质，相邻修正之间的平滑过渡无法保证。例如，图 6 显示，相邻的 USNO A2.0 图块可能具有明显不同的偏差值，其变化可能超过 0.6 角秒。越过这种边界的小行星观测修正将在相应的坐标中出现人为的跳跃。更细的 HEALPix 镶嵌可以用来减轻这种问题。然而，更细的镶嵌意味着每个图块中用于估计修正幅度的恒星数量减少。为了避免减少每个图块的恒星样本量，我们使用径向基函数 (RBFs) (Schaback, 1995) 将来自相邻图块的数据插值到更细的 HEALPix 镶嵌上。这减少了沿图块边界的人为不连续性，同时在每个模板图块中保留了足够数量... (此处文字在 md 文件中不全，我将根据 context 补全) 足够数量的恒星。
%
\begin{figure}
\centering\includegraphics[width=0.7\linewidth]{figs/interp3.pdf} 
\caption{Zoom into a region with a large jump between positive and negative catalog corrections in $\Delta \alpha_{J2000}$ for the USNO A2.0 catalog. Smoothing can avoid quasi-discontinuous transitions between HEALPix tiles with opposite bias values.  \label{fig:smooth}}
\end{figure}
%%%
\section{Ephemeris Prediction Test}
\textbf{【译】星历预报测试}
\label{sec:eph}
%
\begin{figure}
\begin{tabular}{cc}
\includegraphics[width=0.5\linewidth]{figs/1-2_sma.png} &
\includegraphics[width=0.5\linewidth]{figs/4_sma.png}
\end{tabular}
\caption{Cumulative distributions of orbit prediction errors for the debiasing schemes presented in this work (blue line), \citet{farnocchia15} (red line) and the theoretical Gaussian (black line) are shown. The first and second apparition have been used in the left graph and the penultimate apparition in the right graph. \label{fig:eph_pred}}
\end{figure}
%
We follow the example of \citet[][sec. 6]{chesley10} and \citet[][sec. 4.3]{farnocchia15} and evaluate the impact of the new debiasing scheme on asteroid orbit determination and prediction. To achieve that we compare the prediction performance for an ensemble of 700 near-Earth objects (NEOs), each of which had at least five apparitions since 1998. Every one of those apparitions has at least ten observations over 15 days.
Using the astrometric observations from either the first two apparitions or the penultimate apparition, respectively, we construct predictions close to the center of the third apparition, near the middle of the true observational arc. 
The predictions are then compared to the reference solution that uses the full arc. The analysis is performed in a consistent way such that a given debiasing scheme is used for both the short arcs derived from a subset of apparitions and the full arc. All astrometric data is weighted according to the scheme presented in \citet{verevs2017}.

我们遵循 Chesley 等人 (2010) (第 6 节) 和 Farnocchia 等人 (2015) (第 4.3 节) 的示例，评估新去偏方案对小行星轨道确定和预报的影响。为此，我们比较了 700 个近地天体 (NEO) 的预报性能，每个天体自 1998 年以来至少有 5 次出现。每一次出现都至少有 15 天的 10 次观测。使用前两次出现或倒数第二次出现的天体测量观测，我们分别构建接近第三次出现中心（真实观测弧段中间附近）的预报。然后将预报与使用完整弧段的参考解进行比较。分析以一致的方式进行，即给定的去偏方案既用于源自部分出现的短弧，也用于完整弧。所有天体测量数据都根据 Vereš 等人 (2017) 提出的方案进行加权。

Figure \ref{fig:eph_pred} shows the corresponding results.
The prediction errors in the orbital semimajor axes of 700 NEOs derived from astrometry debiased with respect to Gaia DR 2 are only slightly smaller that those using the debiasing scheme presented in \citet{farnocchia15}. These results suggest that most of the statistical bias was de facto eliminated with the previous debiasing approach.

图 7 显示了相应的结果。与其理论高斯分布（黑线）相比，基于 Gaia DR 2 去偏的天体测量数据得出的 700 个 NEO 的轨道半长轴预报误差（蓝线）仅略小于使用 Farnocchia 等人 (2015) 中提出的去偏方案（红线）的结果。这些结果表明，大部分统计偏差实际上已被先前的去偏方法消除了。

%
Using interpolated debiasing tables as discussed in the previous section did not significantly affect the above prediction statistics, either. Given the larger size of the interpolated table and the potential loss in computational efficiency we recommend using interpolated values only when quasi-discontinuities in the former debiasing table are an obstacle to high-fidelity orbit determination.  

使用上一节讨论的插值去偏表并没有显著影响上述预报统计数据。考虑到插值表较大的尺寸和计算效率的潜在损失，我们建议仅在原去偏表中的准不连续性成为高保真轨道确定的障碍时才使用插值值。

\section{Conclusions}
\textbf{【译】结论}
\label{sec:conclusions}
The Gaia astrometric catalog published with the second Gaia data release has become the new standard for minor planet astrometry. 
Its uniformity and high quality in both stellar positions and proper motion allow us to improve prior astrometric observations of minor planets by correcting for potential systematics in other astrometric catalogs. 

随着第二次 Gaia 数据发布的 Gaia 天体测量星表已成为小行星天体测量的新标准。其在恒星位置和自行方面的一致性和高质量使我们能够通过修正其他天体测量星表中潜在的系统误差来改进小行星以前的天体测量观测。

In this work we have used Gaia DR 2 as a reference to calculate debiasing tables for 26 astrometric catalogs. New tables for six catalogs (PPM XL, SST RC5, Gaia DR 1, CMC 15, URAT, UCAC 5 and SDSS 8) that had not been studied in \citet{farnocchia15} are included in this work.

在这项工作中，我们使用 Gaia DR 2 作为参考，计算了 26 个天体测量星表的去偏表。本工作包括了 Farnocchia 等人 (2015) 中未研究的六个星表（PPM XL, SST RC5, Gaia DR 1, CMC 15, URAT, UCAC 5 和 SDSS 8）的新表。

The quality of catalog corrections presented here is significantly improved compared to previous work, in particular for the UCAC series.  
The switch to the new debiasing scheme based on Gaia DR 2 delivers only incremental improvement of the quality of near-Earth asteroid orbits over \citet{farnocchia15}, however.

与以前的工作相比，此处提出的星表修正质量得到了显著提高，特别是对于 UCAC 系列。然而，基于 Gaia DR 2 的新去偏方案对近地小行星轨道质量的提升仅比 Farnocchia 等人 (2015) 有微小的改进。

This suggests that further debiasing attempts may not be necessary. 
The procedure suggested in this work can help mitigate some of the systematics introduced by the use of star catalogs other than Gaia. We would like to emphasize that 
the tables are meant as an a posteriori correction to existing astrometric data. They should not be used to debias new astrometric measurments.
We strongly recommend contemporary Solar System astrometry be conducted with Gaia DR 2 and its successors. 

这表明可能不需要进一步的去偏尝试。本工作中建议的程序可以帮助减轻因使用 Gaia 以外的星表而引入的一些系统误差。我们要强调的是，这些表格旨在作为对现有天体测量数据的后验修正。它们不应用于对新的天体测量结果进行去偏。我们强烈建议当代的太阳系天体测量应使用 Gaia DR 2 及其后续版本进行。

\section{Data Availability}
Tables containing standard resolution and interpolated catalog corrections derived in this work are publicly available at \url{ftp://ssd.jpl.nasa.gov/pub/ssd/debias/debias_2018.tgz}
and 
\url{ftp://ssd.jpl.nasa.gov/pub/ssd/debias/debias_hires2018.tgz}.

\section*{Acknowledgments}
 This research has received funding from the Jet Propulsion Laboratory through the California Institute of Technology postdoctoral fellowship program, under a contract with the National Aeronautics and Space Administration.
 VizieR catalogue access tools, CDS, Strasbourg, France, have been essential to this research project \citep{vizier}.
 This work made use of the IAU Minor Planet Center observations database.
%%%\section*{References}

\bibliography{refs}

%%%
 \section*{Appendix}
 \textbf{【译】附录}
 The appendix shows local position and proper motion systematics for stellar catalogs not presented in the main body of the article (Figures \ref{fig:2mass}-\ref{fig:urat1}). For catalogs without proper motion only the bias values in right ascension and declination are provided. Proper motion corrections for those catalogs resemble those of USNO A2.0 displayed in Figure \ref{fig:usnoa2}.

 附录显示了正文中未展示的恒星表（图 16 到 24）的局部位置和自行系统误差。对于没有自行的星表，仅提供赤经和赤纬的偏差值。这些星表的自行修正类似于图 1 中显示的 USNO A2.0 的自行修正。
 \begin{figure}
 \begin{tabular}{ll}
 \includegraphics[width=0.5\linewidth]{figs/2mass_0.png} &
 \includegraphics[width=0.5\linewidth]{figs/2mass_1.png} \\
 \includegraphics[width=0.5\linewidth]{figs/cmc14_0.png} &
 \includegraphics[width=0.5\linewidth]{figs/cmc14_1.png} \\
 \includegraphics[width=0.5\linewidth]{figs/cmc15_0.png} &
 \includegraphics[width=0.5\linewidth]{figs/cmc15_1.png} \\
 \end{tabular}
 \caption{Corrections in stellar positions for the 2MASS catalog (top panels), CMC 14
 (center panels) and CMC 15 (bottom panels). None of the above catalogs come with proper motion.\label{fig:2mass}}
 \end{figure}

 \begin{figure}
 \begin{tabular}{ll}
 \includegraphics[width=0.5\linewidth]{figs/gsc_1_1_0.png} &
 \includegraphics[width=0.5\linewidth]{figs/gsc_1_1_1.png} \\
 \includegraphics[width=0.5\linewidth]{figs/gsc_1_2_0.png} &
 \includegraphics[width=0.5\linewidth]{figs/gsc_1_2_1.png} \\
 \includegraphics[width=0.5\linewidth]{figs/gsc_act_0.png} &
 \includegraphics[width=0.5\linewidth]{figs/gsc_act_1.png} \\
 \end{tabular}
 \caption{Same as Figure \ref{fig:2mass} for the GSC 1.1 catalog (top panels), the GSC 1.2 catalog
 (center panels) and GSC ACT (bottom panels).  \label{fig:gsc}}
 \end{figure}

 \begin{figure}
 \begin{tabular}{ll}
 \includegraphics[width=0.5\linewidth]{figs/usno_a1_0.png} &
 \includegraphics[width=0.5\linewidth]{figs/usno_a1_1.png} \\
 \includegraphics[width=0.5\linewidth]{figs/usno_sa1_0.png} &
 \includegraphics[width=0.5\linewidth]{figs/usno_sa1_1.png} \\
 \includegraphics[width=0.5\linewidth]{figs/usno_sa2_0.png} &
 \includegraphics[width=0.5\linewidth]{figs/usno_sa2_1.png} \\
 \end{tabular}
 \caption{Same as Figure \ref{fig:2mass} for the USNO A1.0 (top panels), the USNO SA1.0,
 (center panels) and the USNO SA2.0 catalogs (bottom panels).  \label{fig:usnoa}}
 \end{figure}

 \begin{figure}
 \begin{tabular}{ll}
 \includegraphics[width=0.5\linewidth]{figs/act_0.png} &
 \includegraphics[width=0.5\linewidth]{figs/act_1.png} \\
 \includegraphics[width=0.5\linewidth]{figs/act_2.png} &
 \includegraphics[width=0.5\linewidth]{figs/act_3.png}
 \end{tabular}
 \caption{Corrections in stellar positions and proper motion for the ACT catalog. \label{fig:act}}
 \end{figure}

 \begin{figure}
 \begin{tabular}{ll}
 \includegraphics[width=0.5\linewidth]{figs/gaiadr1_0.png} &
 \includegraphics[width=0.5\linewidth]{figs/gaiadr1_1.png} \\
 \includegraphics[width=0.5\linewidth]{figs/gaiadr1_2.png} &
 \includegraphics[width=0.5\linewidth]{figs/gaiadr1_3.png}
 \end{tabular}
 \caption{Corrections in stellar positions and proper motion for the Gaia DR 1 catalog at epoch J2000. \label{fig:gaiadr1}}
 \end{figure}

% \begin{figure}
% \begin{tabular}{ll}
% \includegraphics[width=0.5\linewidth]{figs/hipparcos_0.png} &
% \includegraphics[width=0.5\linewidth]{figs/hipparcos_1.png} \\
% \includegraphics[width=0.5\linewidth]{figs/hipparcos_2.png} &
% \includegraphics[width=0.5\linewidth]{figs/hipparcos_3.png}
% \end{tabular}
% \caption{Corrections in stellar positions and proper motion for the Hipparcos catalog. %\label{fig:hip}}
% \end{figure}

 \begin{figure}
 \begin{tabular}{ll}
 \includegraphics[width=0.5\linewidth]{figs/nomad_0.png} &
 \includegraphics[width=0.5\linewidth]{figs/nomad_1.png} \\
 \includegraphics[width=0.5\linewidth]{figs/nomad_2.png} &
 \includegraphics[width=0.5\linewidth]{figs/nomad_3.png}
 \end{tabular}
 \caption{Corrections in stellar positions and proper motion for the NOMAD catalog. \label{fig:nomad}}
 \end{figure}

 \begin{figure}
 \begin{tabular}{ll}
 \includegraphics[width=0.5\linewidth]{figs/ppm_0.png} &
 \includegraphics[width=0.5\linewidth]{figs/ppm_1.png} \\
 \includegraphics[width=0.5\linewidth]{figs/ppm_2.png} &
 \includegraphics[width=0.5\linewidth]{figs/ppm_3.png}
 \end{tabular}
 \caption{Corrections in stellar positions and proper motion for the PPM catalog. \label{fig:ppm}}
 \end{figure}

 \begin{figure}
 \begin{tabular}{ll}
 \includegraphics[width=0.5\linewidth]{figs/sdss7_0.png} &
 \includegraphics[width=0.5\linewidth]{figs/sdss7_1.png} \\
 \includegraphics[width=0.5\linewidth]{figs/sdss7_2.png} &
 \includegraphics[width=0.5\linewidth]{figs/sdss7_3.png}
 \end{tabular}
 \caption{Corrections in stellar positions and proper motion for the SDSS DR 7 catalog. \label{fig:sdss7}}
 \end{figure}

 \begin{figure}
 \begin{tabular}{ll}
 \includegraphics[width=0.5\linewidth]{figs/sdss8_0.png} &
 \includegraphics[width=0.5\linewidth]{figs/sdss8_1.png} \\
 \includegraphics[width=0.5\linewidth]{figs/sdss8_2.png} &
 \includegraphics[width=0.5\linewidth]{figs/sdss8_3.png}
 \end{tabular}
 \caption{Corrections in stellar positions and proper motion for the SDSS DR 8 catalog. \label{fig:sdss8}}
 \end{figure}

 \begin{figure}
 \begin{tabular}{ll}
 \includegraphics[width=0.5\linewidth]{figs/sstrc5_0.png} &
 \includegraphics[width=0.5\linewidth]{figs/sstrc5_1.png} \\
 \includegraphics[width=0.5\linewidth]{figs/sstrc5_2.png} &
 \includegraphics[width=0.5\linewidth]{figs/sstrc5_3.png}
 \end{tabular}
 \caption{Corrections in stellar positions and proper motion for the SST RC 5 catalog. \label{fig:sstrc5}}
 \end{figure}

 \begin{figure}
 \begin{tabular}{ll}
 \includegraphics[width=0.5\linewidth]{figs/tycho2_0.png} &
 \includegraphics[width=0.5\linewidth]{figs/tycho2_1.png} \\
 \includegraphics[width=0.5\linewidth]{figs/tycho2_2.png} &
 \includegraphics[width=0.5\linewidth]{figs/tycho2_3.png}
 \end{tabular}
 \caption{Corrections in stellar positions and proper motion for the Tycho 2 catalog. \label{fig:tycho2}}
 \end{figure}

 \begin{figure}
 \begin{tabular}{ll}
 \includegraphics[width=0.5\linewidth]{figs/ucac1_0.png} &
 \includegraphics[width=0.5\linewidth]{figs/ucac1_1.png} \\
 \includegraphics[width=0.5\linewidth]{figs/ucac1_2.png} &
 \includegraphics[width=0.5\linewidth]{figs/ucac1_3.png}
 \end{tabular}
 \caption{Corrections in stellar positions and proper motion for the UCAC 1 catalog. \label{fig:ucac1}}
 \end{figure}

 \begin{figure}
 \begin{tabular}{ll}
 \includegraphics[width=0.5\linewidth]{figs/ucac2_0.png} &
 \includegraphics[width=0.5\linewidth]{figs/ucac2_1.png} \\
 \includegraphics[width=0.5\linewidth]{figs/ucac2_2.png} &
 \includegraphics[width=0.5\linewidth]{figs/ucac2_3.png}
 \end{tabular}
 \caption{Corrections in stellar positions and proper motion for  the UCAC 2 catalog. \label{fig:ucac2}}
 \end{figure}

 \begin{figure}
 \begin{tabular}{ll}
 \includegraphics[width=0.5\linewidth]{figs/ucac3_0.png} &
 \includegraphics[width=0.5\linewidth]{figs/ucac3_1.png} \\
 \includegraphics[width=0.5\linewidth]{figs/ucac3_2.png} &
 \includegraphics[width=0.5\linewidth]{figs/ucac3_3.png}
 \end{tabular}
 \caption{Corrections in stellar positions and proper motion for the UCAC 3 catalog. \label{fig:ucac3}}
 \end{figure}


 \begin{figure}
 \begin{tabular}{ll}
 \includegraphics[width=0.5\linewidth]{figs/ucac5_0.png} &
 \includegraphics[width=0.5\linewidth]{figs/ucac5_1.png} \\
 \includegraphics[width=0.5\linewidth]{figs/ucac5_2.png} &
 \includegraphics[width=0.5\linewidth]{figs/ucac5_3.png}
 \end{tabular}
 \caption{Corrections in stellar positions and proper motion for the UCAC 5 catalog. \label{fig:ucac5}}
 \end{figure}


 \begin{figure}
 \begin{tabular}{ll}
 \includegraphics[width=0.5\linewidth]{figs/urat1_0.png} &
 \includegraphics[width=0.5\linewidth]{figs/urat1_1.png} \\
 \includegraphics[width=0.5\linewidth]{figs/urat1_2.png} &
 \includegraphics[width=0.5\linewidth]{figs/urat1_3.png}
 \end{tabular}
 \caption{Corrections in stellar positions and proper motion for the URAT 1 catalog. \label{fig:urat1}}
 \end{figure}



\end{document}